\section{A first-order semantic-refutation method}
\label{sec:formal-practice-03}

\begin{method}[Small-structure search]
To refute a first-order semantic consequence:
\begin{enumerate}
  \item Negate the conclusion at the metalinguistic level by requiring it not to
        be satisfied.
  \item Expand the semantic clauses for the premises and failed conclusion.
  \item Start with a one-element domain unless a premise or failed conclusion
        requires distinct elements.
  \item Introduce witnesses for existential requirements and counterexamples
        for failed universal requirements.
  \item Interpret each relation and function symbol only as far as the expanded
        conditions require.
  \item Complete every partially specified function interpretation to a
        total function on the domain; the remaining values are
        unconstrained, and any choice may be recorded.
  \item Check every premise from the original formulas, not merely the expanded
        notes.
  \item Check that the conclusion fails under the required variable assignment,
        or in the structure when the conclusion is a sentence.
\end{enumerate}
\end{method}

\begin{example}
The sentence \(\forall x\exists y\,R(x,y)\) does not semantically entail
\(\exists y\forall x\,R(x,y)\).  Let the domain be \(\set{0,1}\), and
interpret \(R\) as \(\set{\tuple{0,0},\tuple{1,1}}\).  Every element is
related to itself, so the premise is true.  No single element is related to
both elements, so the conclusion is false.
\end{example}

